% Գլուխ 2. Մեթոդներ (սևագիր v1)

\chapter{Մեթոդներ}

Այս գլխում ներկայացվում են գրաֆի կապակցվածության որոշման և չափման ալգորիթմները, ինչպես նաև ցանցի հուսալիության հաշվարկի մեթոդները: Բոլոր ալգորիթմները իրականացված են գործնական մասում (Գլուխ 3):

\section{Կապակցվածության ալգորիթմներ}

\subsection{Լայնությամբ և խորությամբ որոնում}

Կապակցվածության բաղադրիչները գտնելու հիմնական գործիքներն են \emph{լայնությամբ որոնումը} (BFS) և \emph{խորությամբ որոնումը} (DFS): Երկուսն էլ սկսում են մեկնարկային $s$ գագաթից և այցելում են նրանից հասանելի բոլոր գագաթները: BFS-ը օգտագործում է հերթ (գագաթները մշակվում են հայտնաբերման կարգով), DFS-ը՝ կույտ (stack):

\begin{center}
\fbox{\begin{minipage}{0.9\textwidth}
\textbf{Ալգորիթմ BFS-Բաղադրիչներ}$(G)$\\
1.\ բոլոր գագաթները նշել «չայցելված»\\
2.\ ամեն չայցելված $s$ գագաթի համար.\\
\hspace*{1.5em}3.\ նոր բաղադրիչ $K=\{s\}$, հերթ $Q=[s]$, նշել $s$-ը\\
\hspace*{1.5em}4.\ քանի դեռ $Q$-ն դատարկ չէ. հանել $v$-ն, ամեն չայցելված հարևան $w$-ի համար՝ նշել, ավելացնել $K$-ին և $Q$-ին\\
\hspace*{1.5em}5.\ $K$-ն ավելացնել պատասխանին
\end{minipage}}
\end{center}

\begin{theorem}
BFS-ը և DFS-ը գտնում են $G$ գրաֆի կապակցվածության բաղադրիչները $O(n+m)$ ժամանակում (հարևանության ցուցակներով ներկայացման դեպքում):
\end{theorem}

\begin{proof}
\emph{Ճշտություն.} $s$-ից սկսած ալգորիթմը այցելում է ուղիղ $s$-ից հասանելի գագաթները. եթե $w$-ն հասանելի է $s$-ից $s=v_0,v_1,\dots,v_k=w$ ուղիով, ապա ինդուկցիայով ամեն $v_i$ կայցելվի, քանի որ $v_{i-1}$-ը մշակվելիս նրա բոլոր չայցելված հարևանները ավելացվում են: Հակառակը, այցելված ամեն գագաթ հասանելի է $s$-ից (այն ավելացվել է արդեն այցելված գագաթի հարևանի դերում): Հետևաբար $s$-ից այցելված բազմությունը ճիշտ $s$-ի բաղադրիչն է:

\emph{Բարդություն.} Ամեն գագաթ ավելացվում է հերթ/կույտ ճիշտ մեկ անգամ, և նրա հարևանության ցուցակը անցնում ենք մեկ անգամ: Գումարային աշխատանքը $\sum_v(1+\deg v)=n+2m=O(n+m)$:
\end{proof}

Արդյունքում $G$-ն կապակցված է այն և միայն այն դեպքում, երբ ալգորիթմը գտնում է մեկ բաղադրիչ: Այս ստուգումը անընդհատ կիրառվում է հուսալիության հաշվարկում (տե՛ս բաժին~2.3 և 2.4), քանի որ ամեն պատահական վիճակի համար պետք է պարզել՝ ցանցը կապակցված է, թե ոչ:

\subsection{Կամուրջների և կտրող կետերի որոնում (Տարյանի ալգորիթմ)}

Կամուրջները և կտրող կետերը (տե՛ս բաժին~1.2) կարելի է գտնել մեկ DFS-ով $O(n+m)$ ժամանակում \cite{tarjan}: Յուրաքանչյուր $v$ գագաթի համար DFS-ի ընթացքում հաշվվում են երկու թիվ.
\begin{itemize}
  \item $\mathrm{disc}(v)$՝ $v$-ի հայտնաբերման համարը,
  \item $\mathrm{low}(v)$՝ ամենափոքր $\mathrm{disc}$-ը, որին կարելի է հասնել $v$-ի DFS-ենթածառից՝ օգտագործելով ամենաշատը մեկ հետադարձ կող (back edge).
\end{itemize}

\begin{theorem}\label{thm:tarjan}
Դիցուք $v$-ն DFS-ածառում $u$-ի զավակ է: Այդ դեպքում
\begin{enumerate}
  \item $\{u,v\}$ կողը կամուրջ է $\iff \mathrm{low}(v)>\mathrm{disc}(u)$.
  \item Եթե $u$-ն ածառի արմատը չէ, ապա $u$-ն կտրող կետ է $\iff$ գոյություն ունի զավակ $v$՝ $\mathrm{low}(v)\ge\mathrm{disc}(u)$. Արմատը կտրող կետ է $\iff$ այն ունի առնվազն 2 զավակ:
\end{enumerate}
\end{theorem}

\begin{proof}[Ապացույցի գաղափար]
$\mathrm{low}(v)>\mathrm{disc}(u)$ նշանակում է, որ $v$-ի ենթածառից ոչ մի կող չի հասնում $u$-ին կամ նրա նախորդներին, այսինքն՝ $v$-ի ենթածառը մնացած գրաֆին կապված է միայն $\{u,v\}$ կողով, ուստի այդ կողը ցիկլում չէ և, ըստ Թեորեմ~\ref{thm:bridge}-ի, կամուրջ է: Հակառակը, եթե $\mathrm{low}(v)\le\mathrm{disc}(u)$, ապա կա հետադարձ կող, որը $\{u,v\}$-ի հետ կազմում է ցիկլ: Կտրող կետի դեպքը նույնն է, բայց թույլ անհավասարությամբ՝ քանի որ $\mathrm{low}(v)=\mathrm{disc}(u)$ դեպքում ենթածառը կապվում է հենց $u$-ին, և $u$-ն հեռացնելիս այն անջատվում է: Արմատի համար յուրաքանչյուր զավակի ենթածառ իրար հետ կապված չէ, եթե զավակները մեկից ավելի են:
\end{proof}

\section{Մինիմալ կտրվածք և Մենգերի թեորեմի կիրառում}

Մենգերի թեորեմը (Թեորեմ~\ref{thm:menger}) կապակցվածության հաշվարկը հանգեցնում է \emph{առավելագույն հոսքի} խնդրին:

\begin{theorem}[Առավելագույն հոսք -- մինիմալ կտրվածք, Ֆորդ--Ֆալքերսոն]
Ուղղորդված գրաֆում $s$-ից $t$ առավելագույն հոսքի մեծությունը հավասար է $s$-ը և $t$-ն բաժանող կտրվածքների նվազագույն թողունակությանը:
\end{theorem}

\textbf{Կողային կապակցվածություն.} Չուղղորդված գրաֆի ամեն կողը փոխարինում ենք երկու հակադիր ուղղությամբ կողերով՝ թողունակություն~1: Այդ դեպքում $s$-ից $t$ առավելագույն հոսքը հավասար է կողային չհատվող ուղիների առավելագույն թվին, հետևաբար (Մենգեր) նաև $s$-ը և $t$-ն բաժանող կողերի նվազագույն թվին: Քանի որ ցանկացած կտրվածք բաժանում է ֆիքսված $s$ գագաթը ինչ-որ $t\ne s$-ից,
\[
\lambda(G)=\min_{t\ne s}\ \mathrm{maxflow}(s,t),
\]
ինչը պահանջում է $n-1$ հոսքի հաշվարկ:

\textbf{Գագաթային կապակցվածություն.} Ամեն $v$ գագաթ փոխարինվում է $v_{in}\to v_{out}$ կողով՝ թողունակություն 1 (ինչպես նաև $s$-ի և $t$-ի համար անսահման): Ամեն $\{v,w\}$ կող դառնում է $v_{out}\to w_{in}$ և $w_{out}\to v_{in}$՝ անսահման թողունակությամբ: Այդ ցանցի առավելագույն հոսքը հավասար է ներքին գագաթներով չհատվող ուղիների առավելագույն թվին: $\kappa(G)$-ն հավասար է նվազագույնին ըստ բոլոր ոչ հարևան $(s,t)$ զույգերի (լրիվ գրաֆի համար $\kappa=n-1$):

\textbf{Հոսքի հաշվարկ (Էդմոնդս--Կարպ).} Ալգորիթմը բազմիցս գտնում է $s$-ից $t$ ամենակարճ (կողերի թվով) ավելացնող ուղի BFS-ով և նրա երկայնքով հոսք է ավելացնում \cite{edmondskarp}: Բարդությունը $O(VE^2)$ է, իսկ միավոր թողունակությամբ գրաֆներում հոսքի արժեքը ոչ ավելի, քան $\min(\deg s,\deg t)$ է, ուստի գործնականում ալգորիթմը արագ է աշխատում:

\section{Հուսալիության հաշվարկի ճշգրիտ մեթոդներ}

Հիշենք Գլուխ~1-ի մոդելը. ամեն կող աշխատում է $p$ հավանականությամբ, անկախ մյուսներից, իսկ $R(G)=R_{\mathrm{all}}(G)$-ը այն հավանականությունն է, որ աշխատող կողերով ենթագրաֆը կապակցված է:

\subsection{Լրիվ թվարկում}

Ամենապարզ մեթոդը՝ անցնել բոլոր $2^m$ վիճակներով, ամեն մեկի համար BFS/DSU-ով ստուգել կապակցվածությունը և գումարել կապակցված վիճակների հավանականությունները (բանաձևը Գլուխ~1-ում): Ժամանակը $O(2^m\cdot m)$, ուստի մեթոդը կիրառելի է միայն $m\lesssim 20$ կողով գրաֆների համար: Այն օգտագործվում է որպես հիմք՝ մյուս մեթոդների ճշտությունը ստուգելու համար:

\subsection{Գործոնավորում (deletion--contraction)}

\begin{theorem}[Գործոնավորման թեորեմ]\label{thm:factoring}
Ցանկացած $e\in E$ կողի համար
\[
R(G)=p\cdot R(G/e)+(1-p)\cdot R(G-e),
\]
որտեղ $G-e$-ն $e$-ի \emph{հեռացումն} է, իսկ $G/e$-ն՝ $e$-ի \emph{սեղմումը} (նրա ծայրակետերի միավորումը և ստացված օղակների հեռացումը): Բազային դեպքեր՝ եթե գրաֆը մեկ գագաթ ունի, $R=1$. եթե նույնիսկ բոլոր մնացած կողերով գրաֆը կապակցված չէ, $R=0$:
\end{theorem}

\begin{proof}
Լրիվ հավանականության բանաձևով՝ պայմանավորվելով $e$-ի վիճակի վրա: Եթե $e$-ն աշխատում է (հավ.~$p$), ապա $G$-ի ենթագրաֆը կապակցված է $\iff$ մնացած աշխատող կողերով ենթագրաֆը $G/e$-ում կապակցված է (քանի որ աշխատող $e$-ն արդեն միավորում է իր ծայրակետերը): Եթե $e$-ն խափանվել է (հավ.~$1-p$), ապա խնդիրը ճիշտ $G-e$-ի հուսալիությունն է: Կողերի անկախության պատճառով մնացած կողերի վիճակների բաշխումը երկու դեպքում էլ նույնն է:
\end{proof}

Ռեկուրսիան կանգ է առնում վաղ՝ երբ գրաֆը դառնում է ոչ կապակցված (արդեն հնարավոր չէ $R>0$) կամ մեկ գագաթանի: Այս էտման շնորհիվ գործոնավորումը գործնականում շատ ավելի արագ է լրիվ թվարկումից (տե՛ս Գլուխ~3), թեև վատագույն դեպքում դարձյալ էքսպոնենցիալ է: Իրոք, ընդհանուր դեպքում խնդիրը \#P-դժվար է \cite{provanball}, ուստի բազմանդամային ճշգրիտ ալգորիթմ ակնկալել չարժե (եթե $\mathrm{P}\neq\mathrm{NP}$):

\subsection{Կտրվածքների վրա հիմնված գնահատականներ}

Ցանցը կապակցված չէ այն և միայն այն դեպքում, երբ գոյություն ունի գագաթների $S\subset V$, $\emptyset\ne S\ne V$ բաժանում, որի բոլոր $\delta(S)$ կտրվածքի կողերը խափանվել են: Սրանից երկու պարզ գնահատական է ստացվում.

\begin{theorem}\label{thm:bounds}
\[
1-\sum_{\{S,\,V\setminus S\}}(1-p)^{|\delta(S)|}\ \le\ R(G)\ \le\ 1-(1-p)^{\lambda(G)},
\]
որտեղ գումարումը կատարվում է $V$-ի բոլոր չկարգավորված բաժանումներով:
\end{theorem}

\begin{proof}
\emph{Վերին գնահատական.} Ցանցը խափանվում է առնվազն այն դեպքում, երբ խափանվում են ինչ-որ \emph{մինիմալ} կտրվածքի բոլոր $\lambda$ կողերը՝ հավանականությամբ $(1-p)^\lambda$: Հետևաբար $1-R\ge(1-p)^\lambda$: \emph{Ստորին գնահատական.} $1-R=\Pr(\exists S:\text{ բոլոր }\delta(S)\text{ կողերը խափանված են})$, և միավորման հավանականությունը չի գերազանցում հավանականությունների գումարին (Բուլի անհավասարություն):
\end{proof}

Երբ $p$-ն մոտ է 1-ին, ստորին գնահատականը շատ ճշգրիտ է (տե՛ս Գլուխ~3), որովհետև գումարում գերիշխում են փոքր կտրվածքները: Սակայն գումարը ունի $2^{n-1}-1$ անդամ, ուստի այն նույնպես հաշվողական առումով ծանր է մեծ $n$-երի դեպքում: Բացի այդ, վերին գնահատականից երևում է կապակցվածության դերը. նույն $p$-ի դեպքում մեծ $\lambda$-ով գրաֆները կարող են ունենալ ավելի մեծ հուսալիություն:

\section{Monte Carlo սիմուլյացիա}

Երբ ճշգրիտ հաշվարկը անհնար է, $R(G)$-ը գնահատվում է սիմուլյացիայով \cite{metropolis}: Մեկ փորձը՝ յուրաքանչյուր կող անկախ կերպով «աշխատող» է հայտարարվում $p$ հավանականությամբ, ստուգվում է կապակցվածությունը (BFS կամ DSU): Դիցուք $X_i\in\{0,1\}$-ը $i$-րդ փորձի արդյունքն է ($X_i=1$, եթե կապակցված է): Այդ դեպքում $\Pr(X_i=1)=R(G)$, և գնահատականը
\[
\widehat R_N=\frac1N\sum_{i=1}^N X_i.
\]

\begin{theorem}
$\widehat R_N$-ը անշեղ է, $\mathbb E\widehat R_N=R$, իսկ նրա դիսպերսիան $\mathrm{Var}\,\widehat R_N=R(1-R)/N$: Կենտրոնական սահմանային թեորեմով մոտավորապես $95\%$ վստահությամբ
\[
\bigl|\widehat R_N-R\bigr|\le 1{,}96\sqrt{\frac{\widehat R_N(1-\widehat R_N)}{N}}.
\]
\end{theorem}

Սա երկու գործնական հետևանք ունի.
\begin{itemize}
  \item \emph{Սխալը $\propto 1/\sqrt N$}. սխալը 10 անգամ փոքրացնելու համար պետք է 100 անգամ շատ փորձ: Օրինակ, $R\approx0{,}9$ և $\varepsilon=0{,}001$ ճշտության համար պետք է $N\approx 1{,}96^2\cdot0{,}9\cdot0{,}1/10^{-6}\approx 3{,}5\cdot10^5$ փորձ:
  \item \emph{Ժամանակը կախված չէ $2^m$-ից}. մեկ փորձը $O(m)$ է (DSU), ուստի մեթոդը կիրառելի է մեծ գրաֆների համար, որտեղ ճշգրիտ մեթոդները անհնար են:
\end{itemize}

Թերությունը կապված է հազվադեպ իրադարձությունների հետ. երբ $R$-ը շատ մոտ է 1-ին, խափանումը (հավանականությամբ $1-R$) հազվադեպ է, և նրա \emph{հարաբերական} սխալը մեծ է, եթե $N$-ը բավարար չէ: Այդ դեպքերի համար գոյություն ունեն դիսպերսիայի նվազեցման մեթոդներ (importance sampling և այլն), որոնք դուրս են այս աշխատանքի շրջանակներից:

\section{Մեթոդների համեմատական ամփոփում}

\begin{table}[h]
\centering
\small
\begin{tabular}{|p{3.2cm}|p{3.2cm}|p{3.6cm}|p{3.4cm}|}
\hline
\textbf{Մեթոդ} & \textbf{Ժամանակ} & \textbf{Արդյունք} & \textbf{Կիրառելիություն}\\
\hline
BFS / DFS & $O(n+m)$ & կապակցված է թե ոչ, բաղադրիչներ & ցանկացած չափի գրաֆ\\
\hline
Տարյանի ալգորիթմ & $O(n+m)$ & կամուրջներ, կտրող կետեր & ցանկացած չափի գրաֆ\\
\hline
Մաքս.\ հոսք (Էդմոնդս--Կարպ) & $O(VE^2)$ & $\lambda$, $\kappa$ & միջին չափի գրաֆներ\\
\hline
Լրիվ թվարկում & $O(2^m m)$ & ճշգրիտ $R$ & $m\lesssim 20$\\
\hline
Գործոնավորում & $\le O(2^m)$, էտումով ավելի արագ & ճշգրիտ $R$ & $m\lesssim 25$\\
\hline
Կտրվածքային գնահատականներ & $O(2^n)$ & $R$-ի ստորին/վերին սահման & փոքր $n$\\
\hline
Monte Carlo & $O(Nm)$ & $R$-ի գնահատական $\pm\,1/\sqrt N$ & մեծ գրաֆներ\\
\hline
\end{tabular}
\caption{Մեթոդների համեմատություն (տեսական)}
\end{table}
